% Chapter 4, Section 1 _Linear Algebra_ Jim Hefferon
%  http://joshua.smcvt.edu/linearalgebra
%  2001-Jun-12
\chapter{相似}
我们已经证明，对于任意同态，都存在基 $B$ 与~$D$，使得表示该映射的矩阵具有分块部分单位形。
\begin{equation*}
  \rep{h}{B,D}
  =
    \begin{pmat}{c|c}
       \text{\textit{单位}}  &\text{\textit{零}}   \\
       \hline
       \text{\textit{零}}      &\text{\textit{零}}
    \end{pmat}
\end{equation*}
该表示把映射描述为将
\( c_1\vec{\beta}_1+\dots+c_n\vec{\beta}_n \) 映到
\(c_1\vec{\delta}_1+\dots+c_k\vec{\delta}_k+\zero+\dots+\zero \)，
其中 $n$ 是定义域的维数，而 \( k \) 是值域的维数。
在这种表示下，
由于矩阵的大部分元素都是零，映射的作用很容易理解。

本章考虑定义域与陪域相同的特殊情形。
此时我们自然要求定义域的基与陪域的基相同。
也就是说，
我们想要一个基 \( B \)，使得 \( \rep{t}{B,B} \) 尽可能简单，
这里“简单”指的是矩阵中含有许多零。
我们会发现，
并不总能得到具有上述分块部分单位形的矩阵，
但我们将发展出一种接近的形式，
即一个近乎对角的表示。













\section{复向量空间}
\index{vector space!over complex numbers}
%<*ReasonForShiftToComplexNumbers>
本章需要我们对多项式作因式分解。
但许多多项式在实数范围内不能分解；
例如，
\( x^2+1 \) 不能分解为两个实系数一次多项式的乘积；
要分解它就需要复数 $x^2+1=(x-i)(x+i)$。

因此在本章中，
我们将使用复数作为标量，
包括向量与矩阵中的元素。
也就是说，我们从研究实数上的向量空间
转向研究
复数上的向量空间。\index{complex numbers!vector space over}
任何实数都是复数，
本章的大多数例子也只用实数，
尽管如此，关键的定理要求标量必须是复数。
%</ReasonForShiftToComplexNumbers>
因此
第一节是对复数的复习。

本书转向取复数为标量这一更一般的框架，
是出于务实的理由：为了继续向前推进，我们必须这样做。
不过，取实数以外的结构中的标量来做向量空间，
是一个有趣而又有用的想法。
从一开始就采用这一方法的精彩论述见
\cite{Halmos} 与 \cite{HoffmanKunze}。






\subsectionoptional{多项式因式分解与复数}
\textit{本小节仅为复习。
  完整的展开，包括证明，
  见 \cite{Ebbinghaus}。}

%<*PolynomialReview0>
考虑多项式\index{polynomial}
$p(x)=c_nx^n+\dots+c_1x+c_0$，其
首项系数\index{polynomial!leading coefficient}
$c_n\neq 0$。
我们称它是一个
\definend{次数}为~$n$\index{polynomial!degree}\index{degree of a polynomial}
的多项式。
若 $n=0$，则 $p$ 是一个
常数多项式\index{polynomial!constant}\index{constant polynomial}
$p(x)=c_0$。
不是零多项式的常数多项式，$c_0\neq 0$，
次数为零。
我们规定零多项式的次数为 $-\infty$。
%</PolynomialReview0>

\begin{remark}
将零多项式的次数定义为 $-\infty$ %, which most authors do,
使得等式 $\text{次数}(fg)=\text{次数}(f)+\text{次数}(g)$
对所有多项式都成立。
\end{remark}

%<*PolynomialReview1>
正如整数有除法运算\Dash 例如“\( 4 \) 除 \( 21 \) 得商 \( 5 \) 余 \( 1 \)”\Dash 多项式也有。
%</PolynomialReview1>

\begin{theorem}[多项式的除法定理] \label{th:EuclidForPolys}
\index{division theorem}\index{polynomial!division theorem}
%<*th:EuclidForPolys>
设 \( p(x) \) 是一个多项式。
若 \( d(x) \) 是非零多项式，则存在 \definend{商}多项式与
\definend{余式}多项式 \( q(x) \) 和 \( r(x) \)，使得
\begin{equation*}
  p(x)=d(x)\cdot q(x)+r(x)
\end{equation*}
其中 \( r(x) \) 的次数严格小于 \( d(x) \) 的次数。
%</th:EuclidForPolys>
\end{theorem}

整数表述
“\( 4 \) 除 \( 21 \) 得商 \( 5 \) 余 \( 1 \)”
的要点在于余数小于 \( 4 \)\Dash 虽然 \( 4 \) 能除 \( 5 \) 次，
却不能除 \( 6 \) 次。
类似地，多项式除法表述中的最后一句也是关键。

\begin{example}
若 \( p(x)=2x^3-3x^2+4x \) 且 \( d(x)=x^2+1 \)，则 \( q(x)=2x-3 \) 且
\( r(x)=2x+3 \)。
注意 \( r(x) \) 的次数比 \( d(x) \) 的次数低。
\end{example}

\begin{corollary} \label{co:PolyDividedByLinearPolyIsConstant}
%<*co:PolyDividedByLinearPolyIsConstant>
\( p(x) \) 除以 \( x-\lambda \) 所得的余式是常数多项式
\( r(x)=p(\lambda) \)。
%</co:PolyDividedByLinearPolyIsConstant>
\end{corollary}

\begin{proof}
%<*pf:PolyDividedByLinearPolyIsConstant>
余式必是常数多项式，
因为它的次数小于除式 \( x-\lambda \) 的次数。
为确定这个常数，取定理中的除式
$d(x)$ 为 $x-\lambda$，并将 \( \lambda\/ \) 代入 $x$。
%</pf:PolyDividedByLinearPolyIsConstant>
\end{proof}

%<*PolynomialFactor>
若除式 \( d(x) \) 整除被除式 \( p(x) \)，
即 \( r(x) \) 是零多项式，
则称 \( d(x) \) 是 \( p(x) \) 的一个
因式\index{factor}\index{polynomial!factor}。
该因式的任何根\index{root}\index{polynomial!root}，即任何满足
\( d(\lambda)=0 \) 的 \( \lambda\in\Re \)，
都是 \( p(x) \) 的根，因为
\( p(\lambda)=d(\lambda)\cdot q(\lambda)=0 \)。
%</PolynomialFactor>

\begin{corollary}  \label{co:RootOfPolyIsAssocLinearFactor}
%<*co:RootOfPolyIsAssocLinearFactor>
若 \( \lambda \) 是多项式 \( p(x) \) 的根，
则 \( x-\lambda \) 整除 \( p(x) \)，也就是说，
$x-\lambda$ 是 $p(x)$ 的因式。
%</co:RootOfPolyIsAssocLinearFactor>
\end{corollary}

\begin{proof}
%<*pf:RootOfPolyIsAssocLinearFactor>
由上面的推论，$p(x)=(x-\lambda)\cdot q(x)+p(\lambda)$。
由于 $\lambda$ 是根，
$p(\lambda)=0$，故 $x-\lambda$ 是因式。
%</pf:RootOfPolyIsAssocLinearFactor>
\end{proof}

多项式的重根是这样的数 $\lambda$：该多项式能被 $(x-\lambda)^n$ 整除，其中 $n$ 是大于一的某个幂。
最大的这样的幂称为~$\lambda$ 的
重数\index{multiplicity, of a root}\index{polynomial!multiplicity}。

求一个高次多项式的根与因式可能很难。
但对二次多项式我们有求根公式：~
\( ax^2+bx+c \) 的根就是下面这些
\begin{equation*}
   \lambda_1=\frac{-b+\sqrt{b^2-4ac}}{2a}
   \qquad
   \lambda_2=\frac{-b-\sqrt{b^2-4ac}}{2a}
\end{equation*}
（若判别式 \( b^2-4ac \) 为负，则该多项式没有实数根）。
不能分解成两个次数更低的实系数多项式之积的多项式，
称为在实数上不可约。\index{irreducible polynomial}\index{polynomial!irreducible}

\begin{theorem}  \label{th:CubicsAndHigherFactor}
%<*th:CubicsAndHigherFactor>
任何常数多项式或一次多项式在实数上都不可约。
二次多项式在实数上不可约当且仅当它的
判别式为负。
没有三次或更高次的多项式在实数上不可约。
%</th:CubicsAndHigherFactor>
\end{theorem}

\begin{corollary}  \label{co:RealPolysFactorIntoLinearsAndQuads}
%<*co:RealPolysFactorIntoLinearsAndQuads>
任何实系数多项式都能分解为一些实系数的一次多项式与
不可约二次多项式的乘积。
该分解是唯一的；任何两个分解都具有相同的因式
与相同的幂次。
%</co:RealPolysFactorIntoLinearsAndQuads>
\end{corollary}

注意与整数的素因数分解作类比。
在两种情形中，唯一性条款都非常有用。

\begin{example}
由唯一性，我们不必把它们乘出来就知道
\( (x+3)^2(x^2+1)^3 \) 不等于
\( (x+3)^4(x^2+x+1)^2 \)。
\end{example}

\begin{example}
由唯一性，若 \( c(x)=m(x)\cdot q(x) \)，其中
\( c(x)=(x-3)^2(x+2)^3 \) 且 \( m(x)=(x-3)(x+2)^2 \)，
则我们知道 \( q(x)=(x-3)(x+2) \)。
\end{example}

%<*ComplexNumbers>
虽然 \( x^2+1 \) 没有实根，因而在实数范围内不能分解，
但如果我们想象一个根\Dash 习惯上记作 \( i \)，
使得 \( i^2+1=0 \)\Dash 那么 \( x^2+1 \) 就能分解成一次式之积
\( (x-i)(x+i) \)。
把这个根 \( i \) 添加到实数中，并使新系统对加法与乘法封闭，
我们就得到了
复数，\index{complex numbers}
\( \C=\set{a+bi\suchthat \text{$a,b\in\R$ 且 $i^2=-1$}} \)。
%</ComplexNumbers>

对于标量 $z\in\C$，设 $z=a+bi$，
我们称 $a$ 为~$z$ 的 \definend{实部}，
而~$b$ 为 \definend{虚部}。
我们常把复数描绘在
\definend{复平面}\index{complex plane}上，
其中 $a$ 画在 \definend{实轴}\index{real axis}上，即水平轴，
而 $b$ 画在
\definend{虚轴}上，\index{imaginary axis}
即竖直轴。
注意该点到原点的距离就是长度
$|a+bi|=\sqrt{a^2+b^2}$。

%<*ComplexOperations>
回忆复数的加法与标量乘法的定义
\begin{equation*}
   (a+bi)\,+\,(c+di)=(a+c)+(b+d)i
  \qquad
  r\cdot(a+bi)=(ra)+(rb)i
\end{equation*}
（因此减法就是 $(a+bi)-(c+di)=(a-c)+(b-d)i$）。
再回忆复数与复数相乘的定义。
\begin{align*}
     (a+bi)(c+di) &=ac+adi+bci+bd(-1)  \\
                  &=(ac-bd)+(ad+bc)i
\end{align*}
%</ComplexOperations>

\begin{example}
例如，
% \( (1-2i)\,+\,(5+4i)=6+2i \) and
\( (2-3i)(4-0.5i)=6.5-13i \)。
\end{example}

在复数范围内，任何二次多项式都能分解成一次式之积。
\begin{equation*}
   ax^2+bx+c=
   a\cdot \big(x-\frac{-b+\sqrt{b^2-4ac}}{2a}\big)
    \cdot \big(x-\frac{-b-\sqrt{b^2-4ac}}{2a}\big)
\end{equation*}

\begin{example}
二次多项式 \( x^2+x+1 \) 在复数范围内可分解为
两个一次多项式的乘积：
\begin{equation*}
   \big(x-\frac{-1+\sqrt{-3}}{2}\big)
   \big(x-\frac{-1-\sqrt{-3}}{2}\big)
   =
   \big(x-(-\frac{1}{2}+\frac{\sqrt{3}}{2}i)\big)
   \big(x-(-\frac{1}{2}-\frac{\sqrt{3}}{2}i)\big)
\end{equation*}
\end{example}

与实数不同，在 \( \C \) 中没有不可约的二次式。
所有多项式都能完全分解成一次式之积。

\begin{theorem}[代数基本定理] \label{th:FundThmAlg}\index{Fundamental Theorem!of Algebra}
\hspace*{0em plus2em}
%<*th:FundThmAlg>
复系数多项式能分解成一些复系数一次多项式之积。
该分解是唯一的。
%</th:FundThmAlg>
\end{theorem}












\subsection{复表示}
有了上面关于复数的定义，
我们在实向量空间中使用过的所有运算
都原封不动地适用于复标量的向量空间。

\begin{example}
矩阵乘法是一样的，尽管计算可能涉及更多的算术运算。
\begin{multline*}
   \begin{mat}
      1+1i  &2-0i  \\
       i    &-2+3i
   \end{mat}
   \begin{mat}
      1+0i  &1-0i  \\
      3i    &-i
   \end{mat}             \\
 \begin{aligned}
  &=\begin{mat}
     (1+1i)\cdot(1+0i)+(2-0i)\cdot(3i)   &(1+1i)\cdot(1-0i)+(2-0i)\cdot(-i) \\
     (i)\cdot(1+0i)+(-2+3i)\cdot(3i)     &(i)\cdot(1-0i)+(-2+3i)\cdot(-i)
  \end{mat}                                                  \\
  &=\begin{mat}
     1+7i  &1-1i  \\
    -9-5i  &3+3i
  \end{mat}
 \end{aligned}
\end{multline*}
\end{example}

我们将把前几章的内容尽可能地原封不动地搬过来。
例如，我们将把这个
\begin{equation*}
   \sequence{\colvec{1+0i \\ 0+0i \\ \vdots \\ 0+0i},
             \dots,
             \colvec{0+0i \\ 0+0i \\ \vdots \\ 1+0i}}
\end{equation*}
称为 \( \C^n \) 作为 $\C$ 上的向量空间的
\definend{标准基}\index{standard basis}\index{basis!standard}%
\index{standard basis!complex number scalars}，
并仍记作 \( \stdbasis_n \)。
另一个例子是，
$\polyspace_n$ 将是系数为复数的
次数为~$n$ 的多项式组成的向量空间。











































